\documentclass[10pt,a4paper]{article}

\usepackage[T1]{fontenc}
\usepackage[utf8]{inputenc}
\usepackage[ngerman]{babel}
\usepackage[a4paper,margin=14mm]{geometry}
\usepackage{lmodern}
\usepackage[table]{xcolor}
\usepackage{tabularx}
\usepackage{array}
\usepackage{microtype}
\usepackage[most]{tcolorbox}

\pagestyle{empty}
\setlength{\parindent}{0pt}
\setlength{\parskip}{0pt}
\renewcommand{\familydefault}{\sfdefault}

\definecolor{Navy}{HTML}{17324D}
\definecolor{Blue}{HTML}{3B6EA5}
\definecolor{Sky}{HTML}{EAF3FA}
\definecolor{Mint}{HTML}{E8F5EF}
\definecolor{Gold}{HTML}{F3B33D}
\definecolor{Ink}{HTML}{1E2935}
\definecolor{Muted}{HTML}{617080}
\definecolor{Line}{HTML}{D8E1E8}
\definecolor{CodeBg}{HTML}{F5F7F9}
\definecolor{Cell}{HTML}{FFF4DB}
\definecolor{GridHead}{HTML}{E9EDF1}
\definecolor{GridLine}{HTML}{AEBBC6}

\color{Ink}

\newcommand{\sectiontitle}[1]{%
  \vspace{2.3mm}
  {\color{Blue}\bfseries\large #1}\par
  \vspace{0.8mm}
  {\color{Blue}\rule{\linewidth}{0.7pt}}\par
  \vspace{1.2mm}
}

\newcommand{\gridhead}[1]{\cellcolor{GridHead}\textcolor{Muted}{\bfseries\small #1}}
\newcommand{\gridrow}[1]{\cellcolor{GridHead}\textcolor{Muted}{\bfseries\small #1}}
\newcommand{\sheetsetup}{%
  \arrayrulecolor{GridLine}\setlength{\arrayrulewidth}{0.4pt}%
  \renewcommand{\arraystretch}{1.25}}

% Hervorgehobenes $-Zeichen in Formeln
\newcommand{\fix}{\textcolor{Blue}{\bfseries\$}}

\begin{document}

\colorbox{Navy}{%
  \parbox{\dimexpr\linewidth-2\fboxsep\relax}{%
    \vspace{3mm}\color{white}
    \begin{tabularx}{\linewidth}{@{}Xr@{}}
      {\small\bfseries INFORMATIK · KLASSE 9} &
      \colorbox{Gold}{\color{Navy}\small\bfseries\strut LÖSUNGEN}
    \end{tabularx}
    \vspace{2.5mm}

    {\fontsize{22}{25}\selectfont\bfseries Sicherungsblatt: Aufgabe 2a–2d}\par
    \vspace{1mm}
    {\large Relative und absolute Zellbezüge}
    \vspace{3mm}
  }%
}

\vspace{3mm}
\colorbox{Sky}{%
  \parbox{\dimexpr\linewidth-2\fboxsep\relax}{%
    \vspace{1.8mm}
    \textcolor{Blue}{\bfseries Ziel:}
    Du nutzt relative, absolute und gemischte Zellbezüge beim Auto-Ausfüllen.
    \vspace{1.8mm}
  }%
}

\sectiontitle{1 · Relative Zellbezüge: wandern mit}

\begin{minipage}[t]{\dimexpr\linewidth-62mm\relax}
  \vspace{2mm}
  Zellbezüge \textbf{ohne \$-Zeichen} passen sich beim Ausfüllen an.\par
  \vspace{1.5mm}
  {\small Aufgabe 2c: Fibonacci}
\end{minipage}\hfill
\begin{minipage}[t]{56mm}
  \vspace{0pt}
  \sheetsetup
  \begin{tabular}{|c|l|}
    \hline
    \cellcolor{GridHead} & \gridhead{B} \\ \hline
    \gridrow{7} & \small 0 \\ \hline
    \gridrow{8} & \small 1 \\ \hline
    \gridrow{9} & \cellcolor{Cell}\ttfamily\small =B7+B8 \\ \hline
    \gridrow{10} & \ttfamily\small =B8+B9 \\ \hline
  \end{tabular}
\end{minipage}

\sectiontitle{2 · Absolute Zellbezüge: bleiben fest}

\begin{minipage}[t]{\dimexpr\linewidth-98mm\relax}
  \vspace{2mm}
  \texttt{\fix D\fix 8} meint immer genau D8.\par
  \vspace{1.5mm}
  {\small Aufgabe 2a: Rechnungs\-formular}
\end{minipage}\hfill
\begin{minipage}[t]{92mm}
  \vspace{0pt}
  \sheetsetup
  \begin{tabular}{|c|l|c|l|}
    \hline
    \cellcolor{GridHead} & \gridhead{C} & \gridhead{D} & \gridhead{E} \\ \hline
    \gridrow{1} & \small Artikelbezeichnung & \small netto & \small brutto \\ \hline
    \gridrow{2} & \small Radiergummi & \small 1,20 & \cellcolor{Cell}\ttfamily\small =D2+D2*\fix D\fix 8 \\ \hline
    \gridrow{3} & \small Füller & \small 12,95 & \ttfamily\small =D3+D3*\fix D\fix 8 \\ \hline
    \gridrow{\dots} & & & \\ \hline
    \gridrow{8} & \small MwSt & \small 19 \% & \\ \hline
  \end{tabular}
\end{minipage}

\vspace{1mm}
\colorbox{Sky}{\parbox{\dimexpr\linewidth-2\fboxsep\relax}{%
  \textcolor{Blue}{\bfseries Typischer Fehler:}
  Aus \texttt{=D2+D2*D8} wird in E3 \texttt{=D3+D3*D9} -- D9 ist leer.
  Ebenso muss in Aufgabe 2b der Stundenlohn fest bleiben: \texttt{=C6*\$C\$3}.
}}

\sectiontitle{3 · Gemischte Zellbezüge}

\begin{minipage}[t]{\dimexpr\linewidth-98mm\relax}
  \vspace{2mm}
  \texttt{\fix A1}: Spalte fest.\par
  \vspace{1mm}
  \texttt{A\fix 1}: Zeile fest.\par
  \vspace{1.5mm}
  {\small Aufgabe 2d: Einmaleins}
\end{minipage}\hfill
\begin{minipage}[t]{92mm}
  \vspace{0pt}
  \sheetsetup
  \begin{tabular}{|c|c|l|l|}
    \hline
    \cellcolor{GridHead} & \gridhead{A} & \gridhead{B} & \gridhead{C} \\ \hline
    \gridrow{1} & & \small 1 & \small 2 \\ \hline
    \gridrow{2} & \small 1 & \cellcolor{Cell}\ttfamily\small =\fix A2*B\fix 1 & \ttfamily\small =\fix A2*C\fix 1 \\ \hline
    \gridrow{3} & \small 2 & \ttfamily\small =\fix A3*B\fix 1 & \ttfamily\small =\fix A3*C\fix 1 \\ \hline
  \end{tabular}
\end{minipage}

\sectiontitle{4 · Überblick: Ausfüllen einer Formel aus B2}

\renewcommand{\arraystretch}{1.15}
\begin{tabularx}{\linewidth}{@{}>{\ttfamily\bfseries}p{26mm}>{\ttfamily}p{34mm}>{\ttfamily}p{34mm}X@{}}
  \rowcolor{Sky}\normalfont\bfseries Formel in B2 & \normalfont\bfseries nach unten (B3) & \normalfont\bfseries nach rechts (C2) & \bfseries Art des Bezugs \\
  =A1 & =A2 & =B1 & relativ \\
  \rowcolor{CodeBg}=\fix A\fix 1 & =\$A\$1 & =\$A\$1 & absolut \\
  =\fix A1 & =\$A2 & =\$A1 & gemischt: Spalte fest \\
  \rowcolor{CodeBg}=A\fix 1 & =A\$1 & =B\$1 & gemischt: Zeile fest \\
\end{tabularx}

\sectiontitle{5 · Merke}

\begin{tcolorbox}[colback=Mint,colframe=Blue!40,arc=3mm,boxrule=0.7pt,
  left=3mm,right=3mm,top=2mm,bottom=2mm]
  Relative Zellbezüge passen sich beim Ausfüllen an, absolute wie
  \texttt{\$D\$8} nicht.\par
  \vspace{1.2mm}
  Werte für alle Zeilen (Steuersatz, Stundenlohn) brauchen einen absoluten
  Bezug; \texttt{\$A1} hält die Spalte fest, \texttt{A\$1} die Zeile.
\end{tcolorbox}

\vfill
{\color{Line}\rule{\linewidth}{0.5pt}}\par
\vspace{1.5mm}
\begin{tabularx}{\linewidth}{@{}Xr@{}}
  {\small\color{Muted}Klasse 9 · Tabellenkalkulation} &
  {\small\color{Muted}Sicherungsblatt · Aufgabe 2a–2d · Version 2}
\end{tabularx}

\end{document}
